\documentclass[11pt,oneside,openany]{book}
\usepackage[a4paper,margin=1in]{geometry}
\usepackage[T1]{fontenc}
\usepackage{lmodern}
\usepackage{microtype}
\usepackage{amsmath,amssymb,amsthm,mathtools}
\usepackage[table]{xcolor}
\usepackage{tikz}
\usetikzlibrary{arrows.meta,positioning,calc,decorations.pathreplacing,shapes.geometric}
\usepackage{tcolorbox}
\tcbuselibrary{breakable}
\usepackage{enumitem}
\usepackage{booktabs}
\usepackage{titlesec}
\usepackage{fancyhdr}
\PassOptionsToPackage{hyphens}{url}
\usepackage[colorlinks=true,linkcolor=blue!50!black,urlcolor=blue!50!black]{hyperref}

\setlength{\parskip}{0.5em}
\setlength{\parindent}{0pt}

% ------------------------------------------------------------------ colours
\definecolor{provedc}{RGB}{30,130,60}
\definecolor{partialc}{RGB}{200,120,0}
\definecolor{openc}{RGB}{190,30,45}
\definecolor{ideac}{RGB}{30,80,170}
\definecolor{pinc}{RGB}{200,30,30}
\definecolor{forbc}{RGB}{120,40,160}
\definecolor{cutc}{RGB}{235,235,235}

% ------------------------------------------------------------------ boxes
\newtcolorbox{plainwords}{breakable,colback=ideac!5,colframe=ideac!70!black,
  fonttitle=\bfseries,title={In plain words}}
\newtcolorbox{keyidea}[1][]{breakable,colback=ideac!8,colframe=ideac,
  fonttitle=\bfseries,title={Key idea#1}}
\newtcolorbox{proved}[1]{breakable,colback=provedc!6,colframe=provedc,
  fonttitle=\bfseries,title={Proved: #1}}
\newtcolorbox{partialbox}[1]{breakable,colback=partialc!7,colframe=partialc,
  fonttitle=\bfseries,title={Partly done: #1}}
\newtcolorbox{openitem}[1]{breakable,colback=openc!5,colframe=openc,
  fonttitle=\bfseries,title={Open: #1}}
\newtcolorbox{meeting}{breakable,colback=yellow!8,colframe=orange!70!black,
  fonttitle=\bfseries,title={For the meeting}}
\newtcolorbox{definitionbox}[1]{breakable,colback=gray!5,colframe=gray!60!black,
  fonttitle=\bfseries,title={#1}}

\newtheorem{theorem}{Theorem}[chapter]
\newtheorem{lemma}[theorem]{Lemma}
\newtheorem{corollary}[theorem]{Corollary}

\newcommand{\T}{\mathbin{\square}}
\newcommand{\Z}{\mathbb{Z}}

\titleformat{\chapter}[display]{\bfseries\Huge\color{ideac!80!black}}
  {\large\color{gray}Chapter \thechapter}{0.4em}{}
\pagestyle{fancy}
\fancyhf{}
\fancyhead[L]{\small\leftmark}
\fancyhead[R]{\small\thepage}
\renewcommand{\headrulewidth}{0.4pt}

\begin{document}

% ================================================================== title
\begin{titlepage}
\centering
\vspace*{2cm}
{\color{ideac!80!black}\Huge\bfseries List Colouring the Torus\par}
\vspace{0.6cm}
{\Large Odd cycle times odd cycle: where we are, how we got here,\\ and what is left\par}
\vspace{1.5cm}
\begin{tikzpicture}[scale=0.55]
  \foreach \i in {0,...,4}{\foreach \j in {0,...,6}{
     \pgfmathtruncatemacro{\c}{mod(\i+\j,3)}
     \ifcase\c \def\fc{red!60}\or \def\fc{blue!50}\or \def\fc{green!50!black!60}\fi
     \node[circle,draw,minimum size=9pt,inner sep=0pt,fill=\fc] (v\i\j) at (\j,-\i) {};
  }}
  \foreach \i in {0,...,4}{\foreach \j in {0,...,5}{\pgfmathtruncatemacro{\k}{\j+1}\draw (v\i\j)--(v\i\k);}}
  \foreach \j in {0,...,6}{\foreach \i in {0,...,3}{\pgfmathtruncatemacro{\k}{\i+1}\draw (v\i\j)--(v\k\j);}}
  \foreach \i in {0,...,4}{\draw[dashed,gray] (v\i0)--++(-0.7,0); \draw[dashed,gray] (v\i6)--++(0.7,0);}
  \foreach \j in {0,...,6}{\draw[dashed,gray] (v0\j)--++(0,0.7); \draw[dashed,gray] (v4\j)--++(0,-0.7);}
\end{tikzpicture}
\vspace{1.5cm}

{\large Prepared for the supervisor meeting, Wednesday 28 October 2026, 4:00 PM\par}
\vspace{0.4cm}
{\large Siddharth Narayan Mishra and Amrutanshu\par}
\vfill
{\small Every number in this book has a script in the repository \texttt{\textasciitilde/work/gt} that
reproduces it. Appendix~\ref{app:verify} lists them.}
\end{titlepage}

\tableofcontents

% ================================================================== reading guide
\chapter*{How to read this book}
\addcontentsline{toc}{chapter}{How to read this book}

This book explains one research problem from the start. It assumes you know what a graph, a
cycle and a proper colouring are, and nothing else.

\begin{itemize}[itemsep=0.3em]
\item Chapters 1 and 2 explain the problem and what other people already proved.
\item Chapter 3 explains the tools, in plain words.
\item Chapter 4 is our main finished result: the case where one of the cycles is a triangle.
\item Chapters 5 to 7 are about the next case, $C_5 \T C_n$: what is proved, what is not, and
what we tried this morning.
\item Chapter 8 is the roadmap with dates. Chapter 9 prepares answers for the meeting.
\end{itemize}

Coloured boxes tell you the status of every claim:

\begin{proved}{example}
Green boxes contain things that are fully proved.
\end{proved}
\begin{partialbox}{example}
Orange boxes contain things that are proved only in part, or proved with computer help that
still needs a human-readable version.
\end{partialbox}
\begin{openitem}{example}
Red boxes contain things that are not proved.
\end{openitem}

\begin{meeting}
\textbf{The one-paragraph summary.} We want to show that every grid on a torus made of two odd
cycles can be coloured from any lists of three colours. The usual polynomial method is known to
fail here. We have a complete proof when one cycle is a triangle ($C_3 \T C_n$, all $n$). For
five-cycles ($C_5 \T C_n$) we have proved every ingredient except one statement about small
two-column pieces, which we call the \emph{window statement}. It is checked by computer for
lists of size 3 drawn from a palette of four colours, and survived a targeted search with palettes of five and six colours. This
morning we learned why it is hard: lists that are almost constant beat every local cutting
argument, so the proof must treat them separately. The general odd case needs the same two
ingredients for every odd length, and is still far away.
\end{meeting}

% ================================================================== chapter 1
\chapter{The problem}

\section{List colouring}

In ordinary colouring, every vertex may use any colour. In \emph{list colouring} each vertex $v$
gets its own list $L(v)$ of allowed colours, and we want a proper colouring where each vertex
uses a colour from its own list. A graph is \emph{$k$-choosable} if this is always possible when
every list has $k$ colours, no matter which colours the lists contain. The smallest such $k$ is
the \emph{choice number} $\chi_\ell(G)$.

\begin{center}
\begin{tikzpicture}[v/.style={circle,draw,minimum size=18pt,inner sep=1pt,font=\small}]
  \node[v,fill=red!30] (a) at (0,0) {$a$};
  \node[v,fill=blue!30] (b) at (2.4,0) {$b$};
  \node[v,fill=green!30] (c) at (1.2,1.8) {$c$};
  \draw (a)--(b)--(c)--(a);
  \node[below=2pt of a,font=\small] {$\{\textcolor{red}{\text{red}},\text{blue}\}$};
  \node[below=2pt of b,font=\small] {$\{\text{red},\textcolor{blue}{\text{blue}}\}$};
  \node[above=2pt of c,font=\small] {$\{\text{red},\text{blue},\textcolor{green!50!black}{\text{green}}\}$};
\end{tikzpicture}
\end{center}
\vspace{-0.5em}
{\small\centering A triangle with lists. The colouring shown uses red, blue, green.\par}

List colouring is harder than ordinary colouring. The adversary picks the lists, and lists can be
chosen to block every colouring. A triangle needs 3 colours in any case, so its choice number is
at least 3.

\section{The torus grid}

Take a cycle $C_m$ with vertices $0,1,\dots,m-1$ and a cycle $C_n$. The \emph{toroidal grid}
$C_m \T C_n$ has a vertex $(i,j)$ for each \emph{row} $i$ and \emph{column} $j$. Each vertex is
joined to its two neighbours in its column, $(i \pm 1, j)$, and its two neighbours in its row,
$(i, j \pm 1)$, where indices wrap around. Picture an $m \times n$ grid whose left edge is glued
to its right edge and top edge to its bottom edge: a doughnut. Every vertex has degree 4.

\begin{center}
\begin{tikzpicture}[scale=0.8,v/.style={circle,draw,fill=white,minimum size=7pt,inner sep=0pt}]
  \foreach \i in {0,...,4}{\foreach \j in {0,...,6}{\node[v] (w\i\j) at (\j,-\i) {};}}
  \foreach \i in {0,...,4}{\foreach \j in {0,...,5}{\pgfmathtruncatemacro{\k}{\j+1}\draw (w\i\j)--(w\i\k);}}
  \foreach \j in {0,...,6}{\foreach \i in {0,...,3}{\pgfmathtruncatemacro{\k}{\i+1}\draw (w\i\j)--(w\k\j);}}
  \foreach \i in {0,...,4}{\draw[dashed,gray,-{Stealth}] (w\i6)--++(0.7,0); \draw[dashed,gray,{Stealth}-] (w\i0)--++(-0.7,0);}
  \foreach \j in {0,...,6}{\draw[dashed,gray,-{Stealth}] (w4\j)--++(0,-0.6); \draw[dashed,gray,{Stealth}-] (w0\j)--++(0,0.6);}
  \node[left=14pt,font=\small] at (w20) {rows $0,\dots,4$};
  \node[above=12pt,font=\small] at (w03) {columns $0,\dots,6$};
  \node[right=16pt,font=\small,gray,align=left] at (w26) {dashed edges\\wrap around};
\end{tikzpicture}

{\small $C_5 \T C_7$: five rows, seven columns, edges wrap around in both directions.}
\end{center}

\section{The conjecture}

\begin{definitionbox}{Conjecture (Cai, Wang and Zhu, 2010)}
Every toroidal grid $C_m \T C_n$ is 3-choosable.
\end{definitionbox}

Since every vertex has degree 4, lists of size 5 are always enough (colour greedily). Lists of
size 4 are also known to be enough. The question is whether 3 is enough.

\begin{itemize}
\item If $m$ or $n$ is \textbf{even}, it is true: Li, Shao, Petrov and Gordeev proved it in 2023.
\item If $m$ and $n$ are \textbf{both odd}, it is open. This is our problem.
\end{itemize}

\begin{definitionbox}{Palette: how to read ``with four colours'' in this book}
In this book every list always has exactly \textbf{3} colours. When a computer result says
``with a palette of four colours'' (or ``four-colour windows''), it means the lists are all chosen
from a fixed pool of four colours, say $\{0,1,2,3\}$, so each list is one of $\{0,1,2\}$,
$\{0,1,3\}$, $\{0,2,3\}$, $\{1,2,3\}$. It never means a 4-colouring. A bigger palette gives the
adversary more freedom to choose awkward lists. 3-choosability means no palette of any size
contains a bad choice of lists, which is why the computer alone cannot settle it.
\end{definitionbox}

\begin{plainwords}
We want to show: if you give every vertex of a doughnut grid (both side lengths odd) any three
colours it is allowed to use, you can always colour the grid so that neighbours differ. Nobody
knows how to do this yet, and the standard method is proved not to work.
\end{plainwords}

% ================================================================== chapter 2
\chapter{What was already known}

\section{The polynomial method (Alon and Tarsi)}

Give the vertices variables $x_v$ and multiply $(x_u - x_v)$ over all edges. Call this product
$P$. A colouring is proper exactly when $P \neq 0$ at that point. The \emph{Combinatorial
Nullstellensatz} says: if some monomial $\prod x_v^{t_v}$ of top degree appears in $P$ with a
nonzero coefficient, then the graph can be coloured from any lists with $|L(v)| \ge t_v + 1$.

For a 4-regular graph with $N$ vertices there are $2N$ edges, so we want the monomial
$\prod x_v^2$ (every vertex exponent 2, lists of size 3).

\begin{proved}{the even case (Li, Shao, Petrov, Gordeev 2023)}
If $m$ or $n$ is even, the coefficient of $\prod x_v^2$ in the polynomial of $C_m \T C_n$ is
nonzero, so the torus is 3-choosable. If $m$ and $n$ are both odd, the coefficient is $0$, and in
fact the \emph{Alon--Tarsi number} is 4: the method can never give 3 for odd tori.
\end{proved}

\begin{plainwords}
There is a famous algebra trick that solves the even case in one stroke. For odd tori the trick
provably breaks. So any proof for odd tori must contain at least one step that is not this
algebra. In our proofs that step is a clever choice made by looking at the actual lists.
\end{plainwords}

\section{The smallest case: Galvin}

$C_3 \T C_3$ is the $3 \times 3$ ``rook's graph'': cells of a $3\times 3$ board, adjacent if in the
same row or column. It is the line graph of $K_{3,3}$, and Galvin's famous theorem (the solution of
the Dinitz problem) shows it is 3-choosable. Galvin used the \emph{kernel method}: orient the
edges cleverly so that every induced subgraph has a \emph{kernel} (an independent set that every
outside vertex points into). Then lists of size ``out-degree $+1$'' are enough.

\begin{proved}{$C_3 \T C_3$ (Galvin 1995)}
$C_3 \T C_3$ is 3-choosable.
\end{proved}

That is where the field stood: even tori done, $C_3 \T C_3$ done, every other odd torus open.

% ================================================================== chapter 3
\chapter{Our tools in plain words}

\section{Cutting the torus open}

The torus is hard because it closes up in both directions. If we delete some columns, what is
left is a \emph{strip} $C_m \T P_\ell$: columns are still cycles, but the rows are now paths. A
strip has fewer edges than a torus of the same size, so the algebra has a little room to spare.

\begin{keyidea}[: spare room at the ends]
A strip $C_m \T P_\ell$ has $m\ell$ vertices and $2m\ell - m$ edges, so the algebra has exactly
$m$ ``units'' of spare room. We can spend them at the two end columns:
\begin{itemize}
\item 2 units \textbf{pin} a vertex: its list shrinks to one colour we choose (shown in
\textcolor{pinc}{red}).
\item 1 unit \textbf{forbids} one colour at a vertex: its list shrinks by one
(shown in \textcolor{forbc}{purple}).
\end{itemize}
Whether a particular way of spending the units works depends on a single number, the
\emph{strip coefficient} $c_S$. If $c_S \neq 0$, it works for every choice of lists.
\end{keyidea}

\section{Computing the strip coefficient}

The strip coefficient can be computed column by column: each column contributes a fixed matrix,
the \emph{transfer matrix} $T$, and the coefficient for a strip of length $\ell$ is roughly
``a row vector times $T^{\ell-2}$ times a column vector''. For $m=3$ this matrix is tiny and the
coefficient has a closed formula. For $m = 5$ the matrix is $30 \times 30$, and we use its
eigenvalues: for long strips the biggest eigenvalue dominates, so once its contribution is nonzero
the coefficient can never vanish again.

\begin{plainwords}
We built a calculator for these numbers, proved for $m=3$ that they are never zero when they
should not be, and proved the same for $m = 5$ using the eigenvalues of a $30\times30$ matrix. The
computer does exact arithmetic, so there is no rounding error anywhere.
\end{plainwords}

\section{The plan of every proof in this book}

\begin{center}
\begin{tikzpicture}[box/.style={draw,rounded corners,align=center,minimum width=3.3cm,minimum height=1cm,font=\small}]
  \node[box,fill=ideac!10] (a) {1. Look at the lists.\\Find a good place to cut.};
  \node[box,fill=ideac!10,right=0.6cm of a] (b) {2. Delete a few columns.\\Colour the strip by algebra\\(pin and forbid).};
  \node[box,fill=ideac!10,right=0.6cm of b] (c) {3. Fill the deleted\\columns by hand.};
  \draw[-{Stealth},thick] (a)--(b); \draw[-{Stealth},thick] (b)--(c);
  \node[box,fill=provedc!10,below=0.8cm of b,minimum width=8cm] (d) {If no good place exists, the lists have a special shape:\\colour those directly (for example ``every row has the same list'').};
  \draw[-{Stealth},thick,dashed] (a) |- (d);
\end{tikzpicture}
\end{center}

Step 1 is the non-algebraic step that the polynomial method cannot do on its own.

% ================================================================== chapter 4
\chapter{The finished case: $C_3 \T C_n$}

\begin{proved}{Theorem A}
For every $n \ge 3$, $C_3 \T C_n$ is 3-choosable. So $\chi_\ell(C_3 \T C_n) = 3$.
\end{proved}

For even $n$ this is the known even case, and $n=3$ is Galvin. The new part is odd $n \ge 5$.
The route is Amrutanshu's; the simplification that made it short is in his work plan; we checked
every step, wrote the strip lemma proof ourselves, and verified the finite lemma exhaustively.

\section{Setting up}

Each column of $C_3 \T C_n$ is a triangle. Let $L$ give every vertex 3 colours. There are two cases,
depending on whether the lists change as you walk along a row.

\section{Case 1: some list changes along a row}

Suppose some row $r$ and column $j$ have a colour $\alpha$ that is in $L(r,j)$ but not in
$L(r,j+1)$. Call the columns $C_0 = j$, $C_1$, $C_2$, $C_3$.

\begin{center}
\begin{tikzpicture}[v/.style={circle,draw,minimum size=13pt,inner sep=0pt,font=\scriptsize,fill=white}]
  \fill[cutc] (1.55,0.45) rectangle (3.45,-2.45);
  \node[font=\small,gray] at (2.5,0.75) {deleted};
  \foreach \c/\x in {0/0,1/2,2/3,3/5}{
    \node[font=\small] at (\x,-2.9) {$C_\c$};
    \foreach \r in {0,1,2}{\node[v] (n\c\r) at (\x,-\r) {};}
    \draw (n\c0)--(n\c1)--(n\c2); \draw (n\c0) to[bend left=40] (n\c2);
  }
  \foreach \r in {0,1,2}{\draw (n0\r)--(n1\r)--(n2\r)--(n3\r);}
  \node[v,fill=pinc!40] at (0,0) {$\alpha$};
  \node[v,fill=forbc!30] at (5,-1) {$\not\beta$};
  \node[left=4pt,font=\small] at (n00) {row $r$};
  \node[anchor=west,font=\small] at (5.35,-1) {row $r'$};
  \draw[decorate,decoration={brace,amplitude=6pt,mirror}] (-0.4,-3.3) -- (0.4,-3.3);
  \draw[decorate,decoration={brace,amplitude=6pt,mirror}] (4.6,-3.3) -- (7.2,-3.3);
  \node[font=\small,align=center] at (3.1,-4.0) {strip: from $C_3$ around the torus back to $C_0$\\(odd length $n-2$, coloured by algebra)};
  \node at (6.5,-2) {$\cdots$};
\end{tikzpicture}
\end{center}

\begin{enumerate}[itemsep=0.3em]
\item Delete $C_1$ and $C_2$. The rest is a strip of odd length $n-2$.
\item Colour the strip by algebra, \textcolor{pinc}{pinning} $(r, C_0)$ to $\alpha$ and
\textcolor{forbc}{forbidding} one colour $\beta$ at $(r', C_3)$, for another row $r'$. The strip
lemma (below) says this works.
\item Fill $C_1$ and $C_2$ by hand. Because $\alpha$ is not in $L(r,C_1)$, the vertex
$(r, C_1)$ keeps all 3 of its colours; every other vertex of $C_1, C_2$ keeps at least 2.
\end{enumerate}

\textbf{The strip lemma for $m=3$.} For odd length $\ell$, the coefficient for ``pin one vertex
at one end, forbid at a different row of the other end'' is exactly $\pm 3^{(\ell-1)/2}$, never
zero. We proved this by a short exact calculation with cube roots of unity: each column has only
6 possible states, and the sum factorises.

\textbf{The two-column (prism) lemma.} The deleted part is a \emph{prism}: two triangles
$u_0u_1u_2$ and $v_0v_1v_2$ joined by three rungs. With 3 colours at $u_0$ and at least 2
elsewhere, the fill fails only in two patterns:

\begin{center}
\begin{tikzpicture}[v/.style={circle,draw,minimum size=15pt,inner sep=0pt,font=\scriptsize}]
  \begin{scope}
    \node[font=\small\bfseries] at (1.5,1.2) {Pattern Z};
    \foreach \t/\x/\y in {0/0/0,1/-0.6/-1,2/0.6/-1}{\node[v] (u\t) at (\x,\y) {$u_\t$};}
    \foreach \t/\x/\y in {0/3/0,1/2.4/-1,2/3.6/-1}{\node[v,fill=openc!20] (v\t) at (\x,\y) {$v_\t$};}
    \draw (u0)--(u1)--(u2)--(u0); \draw (v0)--(v1)--(v2)--(v0);
    \draw (u0)--(v0); \draw (u1) to[bend right=30] (v1); \draw (u2) to[bend right=30] (v2);
    \node[font=\scriptsize,align=center] at (1.5,-2.2) {all three $v$-lists are\\the same 2 colours};
  \end{scope}
  \begin{scope}[xshift=6.5cm]
    \node[font=\small\bfseries] at (1.5,1.2) {Pattern W};
    \node[v,fill=yellow!40] (u0) at (0,0) {$u_0$}; \node[v] (u1) at (-0.6,-1) {$u_1$}; \node[v] (u2) at (0.6,-1) {$u_2$};
    \node[v,fill=yellow!40] (v0) at (3,0) {$v_0$}; \node[v] (v1) at (2.4,-1) {$v_1$}; \node[v] (v2) at (3.6,-1) {$v_2$};
    \draw (u0)--(u1)--(u2)--(u0); \draw (v0)--(v1)--(v2)--(v0);
    \draw[very thick,openc] (u0)--(v0); \draw (u1) to[bend right=30] (v1); \draw (u2) to[bend right=30] (v2);
    \node[font=\scriptsize,align=center] at (1.5,-2.2) {$u_1,u_2$ use up $D$, $v_1,v_2$ use up $E$:\\both $u_0$ and $v_0$ are forced to $\zeta$};
  \end{scope}
\end{tikzpicture}
\end{center}

Both patterns need the two $v$-lists in rows $r'$ and $r''$ to be the \emph{same} two colours.
That can only happen in one specific way, and the colour $y'$ that makes it happen can be read off
the lists in advance. We choose $\beta$ to be exactly that colour. Since the strip was coloured
with $\beta$ forbidden at $(r', C_3)$, the bad pattern cannot occur.

\begin{keyidea}[: why this beats the algebra]
$\beta$ is chosen by looking at the lists. That is the non-algebraic step. The algebra alone
(Alon and Tarsi) cannot do this, which is why it fails on odd tori while this proof works.
\end{keyidea}

We verified the two-column lemma by checking every list assignment of the prism up to renaming of
colours: 236 million cases, 8,335 of them not colourable, all with the two $v$-lists equal.

\section{Case 2: no list changes along any row}

Then every row has a constant list $A_r$. Use Galvin on $C_3 \T C_3$ with these row lists to get
three columns $X, Y, Z$ that fit together pairwise. Colour the columns of $C_3 \T C_n$ as
$X, Y, Z, Y, Z, \dots, Y, Z$. Since $n$ is odd, the last column is $Z$, which fits next to $X$.

\begin{center}
\begin{tikzpicture}[c/.style={draw,minimum width=0.7cm,minimum height=1.2cm,font=\small}]
  \foreach \l/\x/\f in {X/0/red!25,Y/1/blue!20,Z/2/green!25,Y/3/blue!20,Z/4/green!25,Y/5/blue!20,Z/6/green!25}{\node[c,fill=\f] at (\x,0) {$\l$};}
  \draw[-{Stealth},gray,thick] (6.4,0.7) to[bend right=25] node[above,font=\scriptsize]{wraps around: $Z$ meets $X$} (-0.4,0.7);
\end{tikzpicture}
\end{center}

\section{A second proof for $n = 3k$ (Amrutanshu)}

When $n$ is a multiple of 3 (and $k$ odd), there is a completely different proof by the kernel
method: colour vertex $(i,j)$ by $(i+j) \bmod 3$, point column edges to the smaller colour and
row edges to the larger. Every vertex gets out-degree 2, and a careful argument shows every
induced subgraph has a kernel. This gives the stronger property of \emph{3-paintability} (the
online version of list colouring). We checked it line by line, and our own counts of
kernel-perfect orientations agree with it. His note also reports an exhaustive search showing
that for $n = 5, 7, 11, 13$ no orientation of this kind is kernel-perfect, so the kernel method
is limited to the family $3 \mid n$.

\begin{meeting}
\textbf{What to say:} ``We have a complete proof for $C_3 \T C_n$ for every $n$. The polynomial
method provably cannot do it; our proof adds one step that looks at the lists to choose a
forbidden colour. Every piece is either proved in the write-up or a finite statement we checked
exhaustively. A separate kernel-method proof covers $n = 3k$ and even gives paintability.''

\textbf{Still to do before submitting:} the hand proof of the two-column lemma (Amrutanshu has
one), and a careful literature check that $C_3 \T C_n$ for odd $n$ is not already in print.
\end{meeting}

% ================================================================== chapter 5
\chapter{Moving to five-cycles: what we have}

For $C_5 \T C_n$ each column is a 5-cycle instead of a triangle. We follow the same plan and ask
which pieces survive.

\begin{center}
\renewcommand{\arraystretch}{1.3}
\begin{tabular}{p{5.3cm}p{3.6cm}p{4.6cm}}
\toprule
Piece & $m = 3$ & $m = 5$ \\
\midrule
Strip lemma (algebra) & \cellcolor{provedc!15}proved, short formula & \cellcolor{provedc!15}proved, computer-assisted \\
Every row constant & \cellcolor{provedc!15}Galvin & \cellcolor{provedc!15}proved (uses Theorem A) \\
Every column constant & \cellcolor{provedc!15}covered by Cases 1 and 2 & \cellcolor{provedc!15}proved (uses Theorem A) \\
One free row & not needed & \cellcolor{provedc!15}proved by hand (new) \\
Filling the deleted columns & \cellcolor{provedc!15}prism lemma & \cellcolor{openc!12}the window statement \\
\bottomrule
\end{tabular}
\end{center}

\section{Rows constant or columns constant}

\begin{proved}{lists constant along rows or along columns (every odd $m$, $n$)}
If every row has a constant list, colour $C_m \T C_3$ with those row lists using Theorem A (it is
$C_3 \T C_m$ turned sideways), take its three columns $X, Y, Z$, and repeat
$X, Y, Z, Y, Z, \dots$ exactly as before. Columns constant is the same argument turned sideways.
\end{proved}

\section{One free row (new)}

\begin{proved}{one free row (every odd $m \ge 5$, every $n$)}
If all rows except one have the same constant list $A$, the torus can be coloured.
\end{proved}

\begin{center}
\begin{tikzpicture}[v/.style={circle,draw,minimum size=14pt,inner sep=0pt,font=\scriptsize}]
  \foreach \j/\p in {0/0,1/1,2/2,3/0,4/1,5/2}{
    \foreach \i/\q in {0/0,1/1,2/2,3/0}{
      \pgfmathtruncatemacro{\col}{mod(\p+\q,3)}
      \ifcase\col \def\fc{red!30}\or \def\fc{blue!25}\or \def\fc{green!30}\fi
      \node[v,fill=\fc] (n\i\j) at (\j*1.1,-\i) {};
    }
    \node[v,fill=yellow!40] (n4\j) at (\j*1.1,-4) {?};
  }
  \foreach \i in {0,...,4}{\foreach \j in {0,...,4}{\pgfmathtruncatemacro{\k}{\j+1}\draw (n\i\j)--(n\i\k);}}
  \foreach \j in {0,...,5}{\foreach \i in {0,...,3}{\pgfmathtruncatemacro{\k}{\i+1}\draw (n\i\j)--(n\k\j);}}
  \node[left=6pt,font=\small,align=right] at (n10) {rows with\\list $A$};
  \node[left=6pt,font=\small] at (n40) {free row};
  \node[right=8pt,font=\scriptsize,align=left] at (n15) {each column reads\\$a,b,c,a$ from the top};
\end{tikzpicture}
\end{center}

\textbf{How.} Colour the rows with list $A$ so that each column reads $a, b, c, a$ (three colours
of $A$ in some order), shifting the order by one step from column to column. Then the two
neighbours of the free row in each column have the \emph{same} colour, so each vertex of the free
row loses at most one colour and keeps at least two. A cycle with lists of size two can be
coloured unless all its lists are the same two colours. If that happens, rotate all the colours
of $A$ by one step; we showed this always breaks the bad situation.

\section{The strip lemma for five-cycles}

\begin{proved}{strip lemma for $m=5$ (computer-assisted)}
For 74 different ways of spending the 5 units of room (one pin plus three forbidden colours,
placed at various rows of the two end columns), the strip coefficient is nonzero for
\emph{every} odd length.
\end{proved}

\textbf{How.} The coefficients for odd lengths satisfy an exact recurrence of order 6, which comes
from the minimal polynomial of the $30 \times 30$ transfer matrix. For each of the 74 ways the
recurrence's largest root dominates from the start, with a nonzero weight, so the coefficient
never vanishes. Checked again by exact computation up to length 41.

\section{Edge deletion and a computer search}

\begin{proved}{edge deletion for $m = 5, 7$ (computer-assisted)}
Remove one edge from $C_5 \T C_n$ or $C_7 \T C_n$: the result is 3-choosable (even with one
list of size 2). Consequence: in any counterexample, every edge has overlapping lists.
\end{proved}

\begin{proved}{no small counterexample (certified computation)}
$C_5 \T C_5$ cannot be made uncolourable by lists of size 3 drawn from 4 colours.
\end{proved}

% ================================================================== chapter 6
\chapter{The missing piece: the window statement}

\section{What a window is}

For $m = 5$, a cut deletes two columns $C_1, C_2$ next to a column $C_0$ where some row's list
changes. The strip colours $C_0$ and $C_3$, but the strip colouring is \emph{not under our
control} except through the pin and the forbidden colours. So we imagine an \emph{adversary}
who colours $C_0$ and $C_3$ in the worst possible way.

\begin{center}
\begin{tikzpicture}[v/.style={circle,draw,minimum size=12pt,inner sep=0pt,font=\scriptsize,fill=white}]
  \fill[cutc] (1.5,0.5) rectangle (3.5,-4.5);
  \node[font=\small,gray] at (2.5,0.85) {window};
  \foreach \c/\x/\lab in {0/0/{$C_0$ (adversary)},1/2/{$C_1$},2/3/{$C_2$},3/5/{$C_3$ (adversary)}}{
    \node[font=\small] at (\x,-5) {\lab};
    \foreach \r in {0,...,4}{\node[v] (m\c\r) at (\x,-\r) {};}
    \foreach \r in {0,...,3}{\pgfmathtruncatemacro{\s}{\r+1}\draw (m\c\r)--(m\c\s);}
    \draw (m\c0) to[bend left=35] (m\c4);
  }
  \foreach \r in {0,...,4}{\draw (m0\r)--(m1\r)--(m2\r)--(m3\r);}
  \node[v,fill=pinc!40] at (0,0) {$\alpha$};
  \foreach \r in {1,3,4}{\node[v,fill=forbc!30] at (5,-\r) {};}
  \node[right=5pt,font=\scriptsize,align=left,forbc] at (5.3,-2.5) {forbidden\\colours};
\end{tikzpicture}
\end{center}

A window is \emph{good} if some way of spending the 5 units (from the 74 that are proved to work)
makes the fill succeed against every adversary. A \emph{cut} is good if its window is good.

\begin{definitionbox}{The window statement (open)}
If every possible cut of a 3-list assignment of $C_5 \T C_n$ is bad, then the lists are constant
along rows, or constant along columns, or have one free row.
\end{definitionbox}

If this is true, then $C_5 \T C_n$ is 3-choosable for every odd $n$: Chapter 5 handles the three
special shapes, and a good cut handles everything else.

\section{What the computer says}

\begin{partialbox}{the window statement}
\begin{itemize}
\item \textbf{Palette of four colours:} completely true for $n = 5$ and $n = 7$. We checked all
$4^{10}$ windows, found the bad ones, and showed that every assignment in which every cut is bad
is row-constant, column-constant or has one free row (600 such assignments for $n=5$, 2,520 for
$n=7$, all of the last kind).
\item \textbf{Palettes of five and six colours:} new bad windows do appear (992 found next to the four-colour
ones), but they still never combine into anything except the three known shapes.
\item \textbf{All palettes:} not provable by computer. Already a palette of five colours would need about 370
CPU-hours, and there is no upper limit on the number of colours.
\end{itemize}
\end{partialbox}

\begin{plainwords}
The computer agrees with the statement everywhere we can look, but ``everywhere we can look'' is
only small palettes. We need an argument that works for a palette of any size.
\end{plainwords}

\section{What we learned along the way}

\begin{itemize}[itemsep=0.3em]
\item Amrutanshu's suggested way of spending the room (all forbidden colours on one side) is not
enough by itself: it fails on 6,904 of the million four-colour windows. Allowing all 74 proved
ways, some of which forbid colours on the pin's side too, cuts this to a few hundred.
\item A cut can be made in two directions at every place where a list changes. Using both
directions removed every remaining problem except the special shapes.
\item With five or more colours, a careless model of the pinned colour $\alpha$ created 23,280
fake ``new obstructions''. Modelling $\alpha$ exactly removed all of them. Every one of those fake
cases was also colourable, so there was never a counterexample.
\end{itemize}

% ================================================================== chapter 7
\chapter{This morning: attacking the window statement}

\section{Attempt 1: a prism lemma for five-cycles}

For $m=3$ the fill failed in only two patterns. We hoped the same for the five-cycle version
($C_5 \T K_2$, ten vertices). Sampling random lists shows it is \textbf{not} so: the failures
come in many shapes, including ones where no two lists are equal. A hand classification would be
a very long case analysis.

\begin{openitem}{a short characterisation of five-cycle fill failures}
Does not exist in the simple form that worked for triangles.
\end{openitem}

\section{Attempt 2: cut three columns instead of two}

\begin{center}
\begin{tikzpicture}[v/.style={circle,draw,minimum size=11pt,inner sep=0pt,fill=white}]
  \fill[cutc] (1.5,0.5) rectangle (4.5,-4.5);
  \foreach \c/\x/\f in {0/0/white,1/2/white,2/3/provedc!25,3/4/white,4/6/white}{
    \foreach \r in {0,...,4}{\node[v,fill=\f] (k\c\r) at (\x,-\r) {};}
    \foreach \r in {0,...,3}{\pgfmathtruncatemacro{\s}{\r+1}\draw (k\c\r)--(k\c\s);}
    \draw (k\c0) to[bend left=35] (k\c4);
  }
  \foreach \r in {0,...,4}{\draw (k0\r)--(k1\r)--(k2\r)--(k3\r)--(k4\r);}
  \node[v,fill=pinc!40] at (0,0) {};
  \node[font=\small] at (0,-5) {$C_0$}; \node[font=\small] at (2,-5) {$C_1$};
  \node[font=\small] at (3,-5) {$C_2$}; \node[font=\small] at (4,-5) {$C_3$}; \node[font=\small] at (6,-5) {$C_4$};
  \node[font=\scriptsize,align=center,provedc!70!black] at (3,0.9) {middle column keeps\\its full lists};
\end{tikzpicture}
\end{center}

If we delete three columns, the middle one is never touched by the adversary, and we get to choose
how to colour the two outer ones so that the middle one works. This is a much more comfortable
position. With only the pin (no forbidden colours at all), \textbf{none of 400 random four-colour
windows and 1 of 400 five-colour-palette windows was bad}, far better than the two-column cut.

\section{The real obstacle: almost-constant lists}

But when the lists are \emph{almost constant} (say every list is $\{0,1,2\}$ except one or two),
the three-column cut fails too: 75 of 135 windows with one changed list, 2,697 of 8,505 with two.
The same was true for the two-column cut. This is not bad luck: when the deleted columns all carry
the same three colours, the adversary always has enough control to force a clash.

\begin{keyidea}[: two regimes]
No local cutting argument can handle almost-constant lists. A proof must split:
\begin{itemize}
\item \textbf{Generic lists} (far from constant): cut, preferably three columns, and fill by hand.
\item \textbf{Almost-constant lists}: colour directly with a global pattern, the way the
one-free-row lemma does.
\end{itemize}
The work is to make ``almost constant'' precise so that both halves can be proved and they cover
everything.
\end{keyidea}

\begin{openitem}{the window statement for all palettes}
Not proved. The attempt today narrowed it to the two-regime plan above.
\end{openitem}

% ================================================================== chapter: update
\chapter{Late morning: a sharper target}

This chapter was added at 9:30 AM on the day of the meeting. It records what the searches found
while the earlier chapters were being read.

\section{Can every cut fail somewhere new?}

We searched directly for an assignment of $C_5 \T C_5$ or $C_5 \T C_7$ (palettes of five or six colours) in
which every cut fails but which is \emph{not} one of the three special shapes. Sixteen independent
searches, about two million steps in total, never found one. Every search got stuck with at least
two places where a cut still works.

\section{What does a bad window look like?}

We first guessed that a bad window must have at least 7 of its 10 lists equal to one set $X$
(``Conjecture N$'$''). An unbiased search \textbf{disproved} this. Here is one of the three bad
windows it found, with $X = \{2,3,4\}$:

\begin{center}
\begin{tikzpicture}[c/.style={draw,minimum width=1.6cm,minimum height=0.6cm,font=\small}]
  \node[font=\small] at (0,0.6) {$C_1$}; \node[font=\small] at (1.8,0.6) {$C_2$};
  \foreach \r/\a/\b/\fa/\fb in {0/{\{0,1,4\}}/{\{0,2,3\}}/white/white,1/X/X/provedc!20/provedc!20,2/X/X/provedc!20/provedc!20,3/X/X/provedc!20/provedc!20,4/{\{1,2,4\}}/{\{1,3,4\}}/white/white}{
    \node[c,fill=\fa] at (0,-0.6*\r) {$\a$}; \node[c,fill=\fb] at (1.8,-0.6*\r) {$\b$};
    \node[font=\scriptsize,left] at (-0.9,-0.6*\r) {row \r};
  }
  \node[font=\scriptsize,right,align=left] at (2.8,-1.2) {a $3 \times 2$ block\\of identical lists};
\end{tikzpicture}
\end{center}

All the bad windows we have, 1,395 of them from every search so far, share one feature:

\begin{partialbox}{Conjecture N$''$ (strong evidence, not proved)}
If a window is bad, then three consecutive rows carry the \emph{same} 3-set $X$ in \emph{both}
deleted columns: a rigid $3 \times 2$ block of identical lists. In the data the block never
consists of the pinned row together with both of its neighbours.
\end{partialbox}

An unbiased search (nine runs, palettes of five to seven colours) looked for a bad window without
such a block and found none; with a palette of five it came close but never got there.

\begin{plainwords}
A cut only ever fails where the lists contain a small rigid block, three rows tall and two columns
wide, in which every list is the same. If we can prove that by hand, the rest of the $m=5$
problem becomes a question about assignments built from such blocks, which is much smaller and
more structured.
\end{plainwords}

\section{A new theorem we can state today}

\begin{proved}{Theorem B (computer-assisted)}
For every odd $n$, the torus $C_5 \T C_n$ can be coloured from any lists of size 3 that are drawn
from a palette of four colours.
\end{proved}

\textbf{Proof outline.}
\begin{enumerate}[itemsep=0.2em]
\item $n=3$ is Theorem A. Row-constant, column-constant and one-free-row assignments are handled by
the lemmas of Chapter 5 (for every palette).
\item Otherwise some list changes along a row, so there is a cut. All $4^{10}$ four-colour windows
were checked: the bad ones (with the real pinned colour) are exactly 184 windows.
\item Only 64 different columns can sit next to a cut. ``Every cut fails'' is then a closed walk
in a finite graph whose states are three consecutive columns (262,144 states, 270,976 edges). We
found every part of the graph that carries a closed walk of odd length (81 parts) and checked that
each one consists of a single special shape (column-constant or one free row). Since walks of every
length are covered, this settles every odd $n$ at once.
\item So if the assignment is not special, some cut is good, and the strip lemma for all odd lengths
plus the fill give a colouring.
\end{enumerate}
The program that judges windows with all 74 shapes was re-checked by an independent Python brute
force on 40 windows (20 good, 20 bad): all 40 agree.

\begin{plainwords}
This is a complete result for infinitely many tori, under one restriction: the lists come from four
colours. Removing that restriction is exactly the window statement.
\end{plainwords}

\section{Rigid rows: the sharpest form of the target}

Call row $t$ of a window \emph{rigid} if its two deleted lists are the same set, $L(t,C_1) =
L(t,C_2)$. In a rigid row the two vertices have almost no room: one of them is forced to copy the
colour the adversary put next to the other. Three rigid rows in a row chain these forced copies
around the cycle, and the adversary can make the chain contradict itself.

\begin{partialbox}{Lemma R (strong evidence, not proved)}
A cut can only fail if at least three rows of its window are rigid.
\end{partialbox}

Evidence: all 4,731 bad (window, $\alpha$) pairs we know have at least three rigid rows (1,659 have
exactly three, 2,852 have four, 220 have five), and sixteen unbiased searches with palettes of five to
eight colours found no bad window with two or fewer.

\textbf{First step towards a proof.} The generic case is ``no rigid row at all''. Amrutanshu's
single cut shape, with a natural choice of forbidden colours, already handles 81\% of the
four-colour windows of this kind; the rest (for example every $C_1$ list $\{0,2,3\}$ and every $C_2$
list $\{1,2,3\}$) need the shapes that also forbid colours on the pinned side. So even the generic
case needs a careful, several-page case analysis. That is Phase 1 of the plan.

\section{A folding trick}

Suppose rows $i-1$ and $i+1$ have the same list in every column. Then fold the torus: glue those
two rows together. The rows other than $i$ become a copy of $C_3 \T C_n$ (for $m = 5$), which
Theorem A colours. Row $i$ then sees the same colour twice in every column, so each of its vertices
loses only one colour, and an odd cycle with lists of size two can almost always be coloured.

\begin{center}
\begin{tikzpicture}[v/.style={circle,draw,minimum size=12pt,inner sep=0pt,font=\scriptsize}]
  \foreach \r/\f in {0/red!25,1/blue!20,2/green!25,3/red!25}{\node[v,fill=\f] (a\r) at (0,-\r*0.8) {};}
  \node[v,fill=yellow!40] (a4) at (0,-3.2) {?};
  \draw (a0)--(a1)--(a2)--(a3)--(a4); \draw (a4) to[bend left=50] (a0);
  \node[font=\scriptsize,right=4pt] at (a0) {row $i+1$}; \node[font=\scriptsize,right=4pt] at (a3) {row $i-1$};
  \node[font=\scriptsize,right=4pt] at (a4) {row $i$};
  \draw[-{Stealth},thick] (1.6,-1.6)--(3,-1.6) node[midway,above,font=\scriptsize]{fold};
  \foreach \r/\f in {0/red!25,1/blue!20,2/green!25}{\node[v,fill=\f] (b\r) at (4,-\r*0.8-0.4) {};}
  \draw (b0)--(b1)--(b2); \draw (b2) to[bend left=50] (b0);
  \node[font=\scriptsize,right=6pt,align=left] at (b1) {a copy of $C_3 \T C_n$:\\Theorem A colours it};
\end{tikzpicture}
\end{center}

\begin{partialbox}{Folding lemma}
This works whenever $C_3 \T C_n$ stays colourable when one vertex has only two colours
(``Lemma S''). Lemma S held in every test (all 393,216 assignments of $C_3 \T C_3$ with four
colours, and 80,000 random ones on larger tori), but is not yet proved: its hard case is again the
almost-constant one. The folding lemma would extend the one-free-row lemma and gives a way to
reduce $m$ to $m-2$ for lists with this kind of symmetry.
\end{partialbox}

\section{Revised plan for Phase 1}

\begin{itemize}
\item Prove Lemma R by hand (a window with at most two rigid rows always has a good cut),
starting with windows that have no rigid row.
\item Classify the assignments in which every cut fails, using that every column next to a cut
contains such blocks. Because the blocks are rigid, this step should reduce to a finite check.
\item Prove Lemma S, which gives the folding lemma.
\end{itemize}

% ================================================================== chapter 8
\chapter{Roadmap and dates}

\section{Where the whole problem stands}

\begin{center}
\renewcommand{\arraystretch}{1.3}
\begin{tabular}{p{3cm}p{4.2cm}p{6.6cm}}
\toprule
Case & Status & What is missing \\
\midrule
$m$ or $n$ even & \cellcolor{provedc!15}known (2023) & nothing \\
$C_3 \T C_n$ & \cellcolor{provedc!15}proved (ours) & hand proof of the prism lemma, literature check \\
$C_5 \T C_n$ & \cellcolor{partialc!15}reduced to one statement & the window statement (two regimes) \\
$C_7 \T C_n$ & \cellcolor{partialc!15}some pieces & strip lemma for $m=7$, window statement \\
all odd $m, n$ & \cellcolor{openc!12}open & strip lemma for every $m$, window statement for every $m$ \\
\bottomrule
\end{tabular}
\end{center}

\section{The plan}

\begin{center}
\begin{tikzpicture}[st/.style={draw,rounded corners,align=left,text width=11.5cm,font=\small,inner sep=5pt}]
  \node[st,fill=provedc!10] (p0) {\textbf{Phase 0, by Fri 30 Oct.} Finish the $C_3 \T C_n$ paper:
    merge the hand proof of the prism lemma, full literature check, agree authorship.};
  \node[st,fill=partialc!10,below=0.3cm of p0] (p1) {\textbf{Phase 1, by Fri 6 Nov.} Prove Lemma R (a cut can only fail when three rows of the
    window are rigid), starting with the case of no rigid row; prove Lemma S (gives the folding
    lemma).};
  \node[st,fill=partialc!10,below=0.3cm of p1] (p2) {\textbf{Phase 2, by Fri 20 Nov.} Generic regime:
    prove that generic windows have a good three-column cut (hand proof with computer-checked case
    splits), and prove the even-length strip lemma for the shapes used (same method as the
    odd-length one, about a day).};
  \node[st,fill=partialc!10,below=0.3cm of p2] (p3) {\textbf{Phase 3, by Fri 27 Nov.} Assemble
    $C_5 \T C_n$, verify everything, write it up. \textbf{Target for $C_5 \T C_n$: end of November.}};
  \node[st,fill=openc!8,below=0.3cm of p3] (p4) {\textbf{Phase 4, December to January.} General odd
    $m$. Needs (a) a strip lemma that works for every $m$ at once (today it is one computation per
    $m$), and (b) the two-regime argument written uniformly in $m$. \textbf{Decision point mid
    December:} if no uniform strip argument is in sight, do $C_7 \T C_n$ by computer and publish
    $m = 3, 5, 7$.};
  \foreach \a/\b in {p0/p1,p1/p2,p2/p3,p3/p4}{\draw[-{Stealth},thick] (\a)--(\b);}
\end{tikzpicture}
\end{center}

\begin{plainwords}
Honest dates: $C_3 \T C_n$ is done now (paper by the end of this week). $C_5 \T C_n$ is
realistic by the end of November, with real risk in Phase 2. The full odd case has no promised
date: the earliest plausible finish is the end of January 2027, and only if two new uniform
arguments are found. Each phase ends with something we can show, even if the next phase stalls.
\end{plainwords}

% ================================================================== chapter 9
\chapter{Questions to expect, and answers}

\begin{description}[style=nextline,itemsep=0.6em]
\item[Is $C_3 \T C_n$ really new?]
Li, Shao, Petrov and Gordeev (2023) state that the odd case is open, and we have found no later
paper that closes it. A full check of everything citing their paper and the recent dynamic survey
is planned for this week.

\item[Why does your proof work when the Alon--Tarsi method provably fails?]
Our proof uses the algebra only on a strip, which has spare room, and then makes one choice by
looking at the lists (the forbidden colour $\beta$). That choice is not algebraic, so the
impossibility result for the torus does not apply.

\item[Why should we trust the computer parts?]
All arithmetic is exact (integers and the cube roots of unity, no rounding). Each result was
reproduced by at least two independent methods. The finite lemma was checked over every list
assignment up to renaming colours. The SAT result comes with a certificate checked by a second
solver. Every number has a script in the repository.

\item[What exactly is missing for $C_5 \T C_n$?]
One statement: when every cut fails, the lists have one of three special shapes. Everything else
is proved. It is verified for a palette of four colours and survived targeted searches with
palettes of five and six. Since this morning we know more: every bad window we have ever found has at least three
rigid rows (Lemma R) and a $3 \times 2$ block of identical lists (Conjecture N$''$), so the missing
proof has a concrete local target.

\item[Why not just compute it?]
The palette is unbounded: lists always have 3 colours, but they can be drawn from any number of
colours. A palette of five alone would take about 370 CPU-hours, and that would still not cover
six or more.

\item[Could the conjecture be false?]
No counterexample has appeared: none for $C_5 \T C_5$ with lists from a palette of four colours (certified), and every
candidate obstruction found by the searches was colourable. A counterexample would have to beat
every cut, which is very restrictive.

\item[What did you do differently from the original plan?]
We proved that the plan's single way of spending the room fails on some windows, found 74 proved
alternatives, added cuts in both directions, proved the one-free-row lemma, proved the strip lemma
for every odd length when $m=5$, and found today that almost-constant lists need a separate,
global argument.

\item[Who did what?]
The $m=3$ route and the kernel proof are Amrutanshu's; the simplification is in his plan. The
verification, the strip lemma proof write-up, the $m=5$ results and the analysis of the window
statement are ours. Credit is shared equally.
\end{description}

% ================================================================== appendix
\appendix
\chapter{Where every number comes from}\label{app:verify}

All paths are inside \texttt{\textasciitilde/work/gt/odd-torus/}.

\begin{center}\small
\renewcommand{\arraystretch}{1.2}
\begin{tabular}{p{5.6cm}p{4.6cm}p{3.6cm}}
\toprule
Claim & Script & Result \\
\midrule
Strip coefficient, $m=3$ & \texttt{ws-W1/strip\_coeff.py}, \texttt{compare.py} & 0 mismatches in 500 \\
Prism lemma, all assignments & \texttt{ws-W9/prism\_sizes.c}, \texttt{prism.c} & 0 exceptions \\
Case 1 extension check & \texttt{ws-W9/ext3.c} & 0 failures \\
One-free-row construction & \texttt{ws-W9/freerow\_check.py} & 0 failures in 20,000 \\
Strip lemma $m=5$, all odd lengths & \texttt{ws-W7/prove\_strip5.py} & 148 cases proved \\
Edge deletion $m=5,7$ & \texttt{ws-W4/} & see \texttt{REPORT.md} \\
No bad lists from a 4-colour palette on $C_5 \T C_5$ & \texttt{ws-W2/certify.py} & UNSAT, certified \\
Windows, plan shape & \texttt{ws-W8/window5.c} & 6,904 bad of 1,048,576 \\
Windows, 74 shapes & \texttt{ws-W8/window5g.c} & 184 bad (exact $\alpha$) \\
Every cut bad, $n=5,7$ & \texttt{ws-W8/global\_q4\_all\_n.py} & 600 and 2,520, all one free row \\
Palettes of five and six & \texttt{ws-W8/global\_alpha.py} & only the known shapes \\
Five-cycle fill failures & \texttt{ws-W10/prism5\_sample.py} & many different shapes \\
Three-column cut & \texttt{ws-W10/cut3.py}, \texttt{cut3\_nearconst.py} & see Chapter 7 \\
Every cut fails, new shapes? & \texttt{ws-W11/evade.c} & none found (16 runs) \\
Conjecture N$'$ & \texttt{ws-W11/hunt.c} & false (3 counterexamples) \\
Conjecture N$''$ & \texttt{ws-W11/hunt5.c} & 1,395 bad windows agree; hunt found none against \\
Lemma S evidence & \texttt{ws-W11/one\_short.py} & 0 failures \\
Lemma R evidence & \texttt{ws-W11/huntR.c} & none against (16 runs) \\
Theorem B, every odd $n$ & \texttt{ws-W8/alln\_q4.py} & 81 odd components, all special \\
Window verdicts re-checked & \texttt{ws-W8/validate\_shapes.py} & 40 of 40 agree \\
No rigid row, single shape & \texttt{ws-W11/r0\_strategy.c} & 47,976 of 248,832 need more shapes \\
\bottomrule
\end{tabular}
\end{center}

\chapter{Glossary}

\begin{description}[itemsep=0.3em]
\item[Choosable] $G$ is $k$-choosable if it can be coloured from any lists of size $k$.
\item[Torus grid $C_m \T C_n$] An $m \times n$ grid with both pairs of opposite sides glued.
\item[Strip] What is left after deleting some columns: columns are cycles, rows are paths.
\item[Pin] Use 2 units of the strip's spare room to force one vertex to a chosen colour.
\item[Forbid] Use 1 unit to stop one vertex from using one chosen colour.
\item[Strip coefficient] The number that decides whether a pin and forbid pattern works.
\item[Transfer matrix] A matrix that computes strip coefficients one column at a time.
\item[Window] The deleted columns of a cut, together with the adversary's choices around them.
\item[Kernel] An independent set that every vertex outside it points into.
\item[Alon--Tarsi number] The best list size the polynomial method can prove. It is 4 for odd tori.
\item[Paintability] A harder, online version of choosability.
\end{description}

\end{document}
